\subsection{Worked example: a numerical SVD and an exact smoother}
\label{sec:d3-example}
Consider six equally spaced observations and order $d=3$.
With the backward-difference convention evaluated at the
right endpoint, the same signed stencil as in
\eqref{eq:difference} is
$\nabla^3 x_{t+3}=x_{t+3}-3x_{t+2}+3x_{t+1}-x_t$.
The full row-rank operator is
\begin{equation}\label{eq:third-order-matrices}
D_3=
\begin{pmatrix}
-1&3&-3&1&0&0\\
0&-1&3&-3&1&0\\
0&0&-1&3&-3&1
\end{pmatrix}.
\end{equation}
Its Gram matrix $Q_3=D_3^\top D_3$ has rank three
and a three-dimensional quadratic-polynomial nullspace,
but neither Gram construction nor an inverse is needed
to compute the reduced SVD $D_3=U\Sigma V^\top$.
Figure~\ref{fig:d3-matrices} gives a concrete
numerical factorization and the exact smoothing operator
at $\lambda=1$.

\begin{figure*}[t]
\centering
\begingroup
\small
\setlength{\arraycolsep}{3pt}
\begin{equation}\label{eq:d3-numerical-svd}
\begin{aligned}
U&\approx
\begin{pmatrix}
-0.5339&-0.7071&0.4637\\
 0.6557& 0.0000&0.7550\\
-0.5339& 0.7071&0.4637
\end{pmatrix},
&
\Sigma&\approx\diag(6.6652,\,3.7417,\,1.2553),\\
V^\top&\approx
\begin{pmatrix}
 0.0801&-0.3387& 0.6155&-0.6155& 0.3387&-0.0801\\
 0.1890&-0.5669& 0.3780& 0.3780&-0.5669& 0.1890\\
-0.3694& 0.5067& 0.3268&-0.3268&-0.5067& 0.3694
\end{pmatrix}.
\end{aligned}
\end{equation}
\begin{equation}\label{eq:d3-numerical-h}
\begin{aligned}
H_1
&=(I_6+D_3^\top D_3)^{-1}
=I_6-V\diag\!\left(\frac{\sigma_j^2}{1+\sigma_j^2}\right)_{j=1}^{3}V^\top\\
&=\frac1{195}
\begin{pmatrix}
171&47&-8&-18&-8&11\\
47&84&59&19&-6&-8\\
-8&59&84&59&19&-18\\
-18&19&59&84&59&-8\\
-8&-6&19&59&84&47\\
11&-8&-18&-8&47&171
\end{pmatrix}.
\end{aligned}
\end{equation}
\endgroup
\caption{Reduced singular-value decomposition of the
$3\times6$ third-difference operator, and its exact
$6\times6$ smoother for $\lambda=1$. The displayed $U$,
$\Sigma$, and $V^\top$ entries are rounded to four
decimal places; the entries of $H_1$ are exact. Singular
vector signs are conventional, not uniquely determined.}
\label{fig:d3-matrices}
\end{figure*}

The inverse is dense even though $Q_3$ has
half-bandwidth three (at most seven nonzero diagonals
in the general $n$ case): for example,
\begin{equation}\label{eq:third-order-distant-leverage}
(H_1)_{1,6}=\frac{11}{195}\neq0.
\end{equation}
This is why a banded factorization of
$A_\lambda=I+\lambda Q_3$ is preferable to
explicitly forming $H_\lambda$ for a single known
penalty, whereas the retained SVD provides reusable
spectral coordinates for many penalties.

To verify the full \emph{nonzero-reference} formula,
take $y=(0,2,0,1,3,2)^\top$, $\mu=1$, and
$\lambda=1$. Directly solving the normal system,
or using the shift-free SVD formula
\eqref{eq:shiftfree}, gives the same result:
\begin{equation}\label{eq:d3-nonzero-affine-check}
\begin{aligned}
\widehat\tau_{1,1}
&=H_1(y+D_3^\top\mathbf1_3)\\
&=y-V\!\left[
w\odot(V^\top y)-(w\oslash\sigma)
\odot(U^\top\mathbf1_3)\right]\\
&=\frac1{195}
(14,238,243,238,308,519)^\top.
\end{aligned}
\end{equation}
Numerically, both full-precision constructions
agree to within $2\times10^{-14}$ in maximum
absolute component error (the rounded SVD values
printed in Figure~\ref{fig:d3-matrices} are for
inspection only). The deterministic
\texttt{tests/test\_d3\_numeric\_svd.py} regression
test computes the factorization afresh and verifies
the exact inverse, its positive spectrum,
and the nonzero-reference reconstruction.
